\begin{lemma}
    Let $A$ be an abelian group.
    Then there exists an injective abelian group $I$
    and an injective homomorphism $A \to I$.
\end{lemma}
\begin{proof}
    Let $S$ be the set of all abelian group homomorphisms
    from $A$ to $\Q/\Z$.
    Let $I$ be $S$-copies of $\Q/\Z$.
    It is clear that $I$ is injective since $\Q/\Z$ is injective.

    We have a homomorphism $\Phi: A \to I$ given by $a \mapsto (\phi(a))_{\phi \in S}$.
    We want to show that $\Phi$ is injective,
    and it suffices to show that for any nonzero $a \in A$,
    there exists a homomorphism $\phi: A \to \Q/\Z$ such that $\phi(a) \neq 0$.

    Let $a \in A$ be nonzero. Then we have a subgroup $\langle a \rangle$ of $A$
    generated by $a$. We have a homomorphism $\psi: \langle a \rangle \to \Q/\Z$ given by $ma \mapsto m/n +\Z$ where $n$ is the order of $a$.
    By the injectivity of $\Q/\Z$, we can extend $\psi$ to a homomorphism $\phi: A \to \Q/\Z$.
    Then we have a non-zero homomorphism $\phi$ with $\phi(a) = \psi(a) = 1/n + \Z \neq 0$.
    It follows that $\Phi$ is injective.
\end{proof}

\begin{proposition}
    Let $R$ be a ring and $M$ be an $R$-module.
    Then there are enough injective $R$-modules.
\end{proposition}
\begin{proof}
    Let $M$ be a left $R$-module.
    Look at the composition of $R$-module homomorphisms $M \overset{\sim}{\to} \Hom_{R}(R, M) \hookrightarrow \Hom_{\Z}(R, M)$
    where $M \overset{\sim}{\to} \Hom_{\Z}(R, M)$ is given by $m \mapsto (r \mapsto rm)$.
    We take an injective module $I$ as in the previous lemma.
    Then we have an injection $\Hom_{\Z}(R, M) \hookrightarrow \Hom_{\Z}(R, I)$
    via $M \hookrightarrow I$.
    
    It is left to show that $I':=\Hom_{\Z}(R, I)$ is injective as an $R$-module.
    Take any injection $A \hookrightarrow B$ of $R$-modules.
    We want to show that $\Hom_{R}(B, I') \twoheadrightarrow \Hom_{R}(A, I')$.
    We observe that $\Hom_{R}(B, I') = \Hom_{R}(B, \Hom_{\Z}(R, I)) \cong \Hom_{\Z}(B, I)$
    via the adjointness of $\Hom$ functors.
    Similarly, we have $\Hom_{R}(A, I') \cong \Hom_{\Z}(A, I)$.
    Since $I$ is injective as an abelian group (as a $\Z$-module), we have $\Hom_{\Z}(B, I) \twoheadrightarrow \Hom_{\Z}(A, I)$.
    It follows that $\Hom_{R}(B, I') \twoheadrightarrow \Hom_{R}(A, I')$ and $I'$ is injective as an $R$-module.
\end{proof}

Now we work with a $G$-module $A$.
\begin{lemma}
    $A^G = \Hom_{\Z[G]}(\Z, A)$.
\end{lemma}
\begin{proof}
    By adjointness, we have $\Hom_{\Z[G]}(\Z, A) \cong \Hom_{\Z}(\Z, A)^G = A^G$.
\end{proof}

\newpage
\section{Derived Categories}
Since we have enough injective $G$-modules, we can consider the derived functor.
We consider an injective resolution of $A$ as a $\Z[G]$-module
\[
    0 \to A \to I^0 \to I^1 \to I^2 \to \cdots.
\]

From the cochain complex,
we have 
\[
\mathrm{Ext}^i_{\Z[G]}(\Z, A) = H^i(\Hom_{\Z[G]}(\Z, I)).
\]
We note that it doesn't depend on the choice of injective resolution of $A$;
hence it is well-defined.

\begin{definition}
    Define the $i$-th cohomology group of $G$ with coefficients in $A$ as
    $H^i(G, A) := \mathrm{Ext}^i_{\Z[G]}(\Z, A)$.
\end{definition}

Here is a corollary of the previous lemma.
\begin{corollary}
    We have $H^0(G, A) = A^G$.
\end{corollary}

How do we actually compute $H^i(G, A)$?
The hardness of the computation depends on the choice of the injective resolution of $A$.
For $\mathrm{Ext}^i_{\Z[G]}(\Z, A)$,
instead of taking an injective resolution of $A$,
we can take a projective resolution of $\Z$ as a $\Z[G]$-module.
One benefit of this approach is that
projective resolutions are often easier to construct than injective resolutions
since projective modules are direct summands of free modules.

Recall the following theorem.
\begin{theorem}
    Let $M$ be a projective $R$-module and $N$ be an $R$-module.
    Then $\mathrm{Ext}^i_R(M, N) = 0$ for all $i > 0$.
\end{theorem}
\begin{proof}
    We may assume that $M$ is a free $R$-module since projective modules are direct summands of free modules.
    Let $M = \bigoplus_{j \in I} R$ be a free $R$-module.
    Then we have
    \(
        \mathrm{Ext}^i_R(M, N) = \mathrm{Ext}^i_R\left(\bigoplus_{j \in I} R, N\right) = 0
    \)
    for all $i > 0$.
\end{proof}

\begin{corollary}
    If $f: A \to B$ is a $G$-invariant map, then there is a canonical map $f^i: H^i(G, A) \to H^i(G, B)$ for all $i \geq 0$.
\end{corollary}

\begin{corollary}
    $H^0(G, A) = A^G$, and $H^1(G, A)$ is the set of equivalence classes of extensions of $\Z$ by $A$ as $\Z[G]$-modules.    
\end{corollary}

\begin{corollary}
    Let $0 \to A' \to A \to A'' \to 0$ be a short exact sequence of $G$-modules.
    Then we have a long exact sequence
    \[
        \cdots \to H^i(G, A') \to H^i(G, A) \to H^i(G, A'') \to H^{i+1}(G, A') \to \cdots.
    \]
\end{corollary}

We are still working with a finite group $G$.
\begin{definition}
    The \textit{norm element} of $\Z[G]$ is defined as the sum of all elements in $G$:
    \[
        N := \sum_{g \in G} g.
    \]
\end{definition}

\begin{remark}
    \noindent\begin{enumerate}
        \item The norm element $N$ is a central element of $\Z[G]$ since we have
        \[gN = g\sum_{h \in G} h = \sum_{h \in G} gh = N\]
        for any $g \in G$. Likewise $Ng = N$. Hence $gN = N = Ng$.
        \item $N$ is a $G$-invariant element of $\Z[G]$ since we have
        \[g \cdot N = g \cdot \sum_{h \in G} h = \sum_{h \in G} gh = N\]
        for any $g \in G$.
    \end{enumerate}
\end{remark}

\begin{theorem}
    Let $G$ be a finite group of order $m$ and $A$ be a $G$-module
    such that $A$ is also an $R$-module where $R$ is any commutative ring with $1/m \in R$.
    Then $N/m$ is a central idempotent of $R[G]$.
    Moreover, $H^i(G, A) = 0$ for all $i > 0$ and $H^0(G, A) = eA$ where $e = N/m$.
\end{theorem}
\begin{proof}
    We have 
    \begin{align*}
        N^2 &= \left(\sum_{g \in G} g\right)N = \sum_{g \in G} gN = \sum_{g \in G} N = mN.
    \end{align*}
    Therefore,
    \[
        (N/m)^2 = N^2/m^2 = mN/m^2 = N/m
    \]
    i.e. $N/m$ is an idempotent of $R[G]$.
    We also have $N/m$ is central since $N$ is central.

    Then we have $R[G] = R[G]e \oplus R[G](1-e)$ where $e = N/m$ is an idempotent of $R[G]$.
    In particular, $R[G]e$ is a projective $R[G]$-module since it is a direct summand of a free module $R[G]$.
    It follows that $\mathrm{Ext}^i_{R[G]}(R[G]e, A) = 0$ for all $i > 0$.

    Since $1/m \in R$, we have that $A$ is naturally a $\Z[1/m][G]$-module.
    Let $P_\bullet \to \Z$ be a resolution of $\Z$ by free $\Z[G]$-modules.
    Localizing at $m$, we obtain $P_\bullet[1/m] \to \Z[1/m]$
    which is also a free resolution of $\Z[1/m]$ by $\Z[1/m][G]$-modules.

    Since $m$ acts invertibly on $A$, the universal property of localization gives
    an isomorphism
    \(
        \Hom_{\Z[1/m][G]}\big(P_i[1/m], A\big) \xrightarrow{\ \sim\ } \Hom_{\Z[G]}(P_i, A),
        \ f \mapsto f|_{P_i},
    \)
    which is natural in $P_i$: every $\Z[G]$-linear map $P_i \to A$ extends uniquely to
    $P_i[1/m]$ by $f(x/m^k) = m^{-k}f(x)$. These isomorphisms commute with the
    differentials, so
    \begin{align*}
        H^i(G,A) &= H^i\big(\Hom_{\Z[G]}(P_\bullet, A)\big)\\
        &\cong H^i\big(\Hom_{\Z[1/m][G]}(P_\bullet[1/m], A)\big)\\
        &= \operatorname{Ext}^i_{\Z[1/m][G]}(\Z[1/m], A).
    \end{align*}

    Applying the previous result with $R := \Z[1/m]$, 
    we are done with $\Z[1/m] = R[G]e = \Z[1/m][G]e$ as $\Z[1/m][G]$-modules.
    Since $e = N/m$, we know $\Z[1/m][G]e = \Z[1/m]e$. By augmentation, we have $\Z[1/m]e \cong \Z[1/m]$ as $\Z[1/m][G]$-modules.
    Therefore, we have $H^i(G, A) = 0$ for all $i > 0$.
    For $i = 0$, we have 
    \begin{align*}
    H^0(G, A) = \Hom_{\Z[1/m][G]}(\Z[1/m], A) = \Hom_{\Z[1/m][G]}(\Z[1/m][G]e, A) = eA.
    \end{align*}
\end{proof}